\chapter{CPU構成と検証用実装}\label{chap:2}

\section{命令経路と共有データポートの調停}\label{sec:2-1}

各コアはPCと32本の32ビットレジスタを持ち、同じプログラムを実行する。
各コアは自身のPCが指す命令を独立に読み出す。
調停器は両コアからデータアクセスの要求を受け、要求が重なるとCore0を優先する。
Core1は競合中にPCとレジスタを保持する。



\begin{figure}[H]

\centering

\begin{tikzpicture}[x=\dimexpr\linewidth/500\relax,y=-1pt,>=stealth]
\draw[<->,line width=0.5pt] (110,48) -- (110,96);
\draw[<->,line width=0.5pt] (389,48) -- (389,96);
\draw[<->,line width=0.5pt] (163,164) -- (163,211);
\draw[<->,line width=0.5pt] (337,164) -- (337,211);
\draw[<->,line width=0.5pt] (249.5,269) -- (249.5,292);
\draw[<->,line width=0.5pt] (195,332) -- (195,343) -- (122,343) -- (122,355);
\draw[<->,line width=0.5pt] (305,332) -- (305,343) -- (377,343) -- (377,355);
\node[align=center,inner sep=2pt,text width=0.91000\linewidth,minimum height=48pt,draw=black!50,fill=black!3] at (249.5,24.0) {\fontsize{10}{13.0}\selectfont 命令ROM　inst\_memory\\同じ配列を二つのPCから独立に読出し};
\node[align=center,inner sep=2pt,text width=0.35400\linewidth,minimum height=68pt,draw=black!50,fill=black!3] at (110.5,130.0) {\fontsize{10}{13.0}\selectfont Core0　ID 0\\PC・32レジスタ\\coreIDはCPU内で応答};
\node[align=center,inner sep=2pt,text width=0.35400\linewidth,minimum height=68pt,draw=black!50,fill=black!3] at (388.5,130.0) {\fontsize{10}{13.0}\selectfont Core1　ID 1\\PC・32レジスタ\\coreIDはCPU内で応答};
\node[align=center,inner sep=2pt,text width=0.51000\linewidth,minimum height=58pt,draw=black!50,fill=black!3] at (249.5,240.0) {\fontsize{10}{13}\selectfont 共有データ調停器\\Core0固定優先　許可した要求を選択\\競合中のCore1へstall};
\node[align=center,inner sep=2pt,text width=0.27800\linewidth,minimum height=40pt,draw=black!50,fill=black!3] at (249.5,312.0) {\fontsize{9}{11.700000000000001}\selectfont アドレスによる選択\\0x1000未満はROM};
\node[align=center,inner sep=2pt,text width=0.40000\linewidth,minimum height=58pt,draw=black!50,fill=black!3] at (122.0,384.0) {\fontsize{9}{11.700000000000001}\selectfont データ読出し用ROM　data\_image\\命令ROMと同じMEM\_FILE\\ROM宛てのstoreは抑止};
\node[align=center,inner sep=2pt,text width=0.40000\linewidth,minimum height=58pt,draw=black!50,fill=black!3] at (377.0,384.0) {\fontsize{9}{11.700000000000001}\selectfont 共有RAM　data\_memory\\4096ワード　16 KiB\\立上り書込み・組合せ読出し};
\node[align=center,inner sep=2pt,text width=0.16400\linewidth,minimum height=22pt] at (165.0,74.0) {\fontsize{9}{11.700000000000001}\selectfont PC0 / 命令0};
\node[align=center,inner sep=2pt,text width=0.16400\linewidth,minimum height=22pt] at (444.0,74.0) {\fontsize{9}{11.700000000000001}\selectfont PC1 / 命令1};
\node[align=center,inner sep=2pt,text width=0.25800\linewidth,minimum height=32pt] at (86.5,190.0) {\fontsize{9}{11.700000000000001}\selectfont 要求・アドレス・書込値\\読出値・stall};
\node[align=center,inner sep=2pt,text width=0.25800\linewidth,minimum height=32pt] at (417.5,190.0) {\fontsize{9}{11.700000000000001}\selectfont 要求・アドレス・書込値\\読出値・stall};
\end{tikzpicture}


\caption{CPUとメモリの接続。双方向矢印は要求と応答の往復を表す。coreIDは各CPU内で返すため、共有データポートを使わない。}\label{fig:1}

\end{figure}

命令用のinst\_memoryは一つの配列を二つのPCで読み出す。
追加したデータ読出し用のdata\_imageは、同じMEM\_FILEを初期値として持つ別インスタンスである。
調停後のアドレスが0x1000未満なら入力ROMの値を返し、それ以外ならRAMへ接続する。



RAMは4,096ワード、計16 KiBである。
読出しは非同期、書込みはクロックの立上りで行う。
このメモリ構成をFPGA資源へどう対応させるかは、今回のRTLシミュレーションとは別の検討事項である。



\section{メモリ配置とコア識別}\label{sec:2-2}

\begin{table}[H]

\centering

\caption{プログラムが使うメモリ領域}

\label{tab:2}

\footnotesize

\setlength{\tabcolsep}{3pt}

\renewcommand{\arraystretch}{1.25}

\begin{tabular}{>{\raggedright\arraybackslash}p{\dimexpr0.284412\linewidth-2\tabcolsep\relax}>{\raggedright\arraybackslash}p{\dimexpr0.338490\linewidth-2\tabcolsep\relax}>{\raggedright\arraybackslash}p{\dimexpr0.377098\linewidth-2\tabcolsep\relax}}\toprule

\textbf{アドレス} & \textbf{内容} & \textbf{アクセスする主体}\\\midrule

0x0000〜0x07FF & プログラム & 両コアの命令フェッチ\\

0x0800〜0x08FF & 入力ROM　最大64要素 & Core0がlwで初期コピー\\

0x1000 & READY & Core0が1を書き、Core1が読む\\

0x1004 / 0x1008 & LEFT\_DONE / RIGHT\_DONE & 各担当コアが1を書く\\

0x100C & COMPLETE & Core0が1を書き、Core1が読む\\

0x1100〜0x11FF & 配列　最大64ワード & 分担ソートと結果の書戻し\\

0x1300〜0x13FF & マージ用作業領域 & Core0\\

0x1FFC / 0x3000 & coreID　値は0または1 & 各CPU内で応答\\

\bottomrule\end{tabular}

\end{table}

0x1000未満の入力ROMをlwで取得できるよう、コード1の選択信号でRAMへのアクセスを抑止する。
ROM宛てのstoreもRAMへ伝えない。
読出し値は同じ選択信号でROMまたはRAMから選ぶ。



\par\medskip\noindent\begin{minipage}{\linewidth}
\par\medskip\noindent{\small\bfseries コード1　ROMの判定とRAMアクセスの抑止　simple\_memory\_arbiter.v}\par\nopagebreak

\lstinputlisting{code/code-01.txt}
\end{minipage}\par\medskip

プログラムは0x1FFCでコアIDを読む。
コード2では、従来の0x3000と同じCPU内部応答へ接続した。
両アドレスへのアクセスは共有データ要求を出さず、storeは無視する。



\par\medskip\noindent\begin{minipage}{\linewidth}
\par\medskip\noindent{\small\bfseries コード2　コアIDの別名とデータ要求　simple\_cpu\_with\_mem\_interface.v}\par\nopagebreak

\lstinputlisting{code/code-02.txt}
\end{minipage}\par\medskip

本文と付録の行番号は参照元の実ファイルに対応する。
離れた範囲を抜粋した箇所は...で示し、本文コードでは説明用コメントを省略した。
命令のPCと機械語を含む一覧は付録Aに置く。


